\title{2025 年高考数学 上海卷}
\date{}
\maketitle

\section*{一、填空题}
本大题共 12 题，第 1~6 题每题 4 分，第 7~12 题每题 5 分，共 54 分.考生应在答题纸的相应位置直接填写结果

\begin{question}[points=4]
已知全集 $U = \{x \mid 2 \leq x \leq 5, x \in \mathbf{R}\}$ ，集合 $A = \{x \mid 2 \leq x < 4, x \in \mathbf{R}\}$ ，则 $\overline{A} =$ ．
\fillin[{$\{x \mid 4 \leq x \leq 5, x \in \mathbf{R}\}$  $[4,5]$}]
\begin{solution}
\textbf{答案：}$\{x \mid 4 \leq x \leq 5, x \in \mathbf{R}\}$  $[4,5]$

\textbf{分析：}根据补集的含义即可得到答案.

\textbf{详解：}根据补集的含义知 $\overline{A} = \{x \mid 4 \leq x \leq 5, x \in \mathbf{R}\}$ .
\end{solution}
\end{question}

\begin{question}[points=4]
不等式 $\frac{x-1}{x-3} < 0$ 的解集为．
\fillin[{$(1,3)$}]
\begin{solution}
\textbf{答案：}$(1,3)$

\textbf{分析：}转化为一元二次不等式 $(x-1)(x-3) < 0$ ，解出即可.

\textbf{详解：}原不等式转化为 $(x-1)(x-3) < 0$ ，解得 $1 < x < 3$ ，则其解集为 $(1,3)$ .
\end{solution}
\end{question}

\begin{question}[points=4]
己知等差数列 $\{a_n\}$ 的首项 $a_1 = -3$ ，公差 $d = 2$ ，则该数列的前 $6$ 项和为．
\fillin[{$12$}]
\begin{solution}
\textbf{答案：}$12$

\textbf{分析：}直接根据等差数列求和公式求解.

\textbf{详解：}根据等差数列的求和公式，$S_6 = 6a_1 + \frac{6 \times 5}{2} d = 12$ .
\end{solution}
\end{question}

\begin{question}[points=4]
在二项式 $(2x-1)^5$ 的展开式中，$x^3$ 的系数为．
\fillin[{$80$}]
\begin{solution}
\textbf{答案：}$80$

\textbf{分析：}利用通项公式求解可得.

\textbf{详解：}由通项公式 $T_{r+1} = \mathrm{C}_5^r \cdot 2^{5-r} \cdot x^{5-r} \cdot (-1)^r = \mathrm{C}_5^r \cdot (-1)^r \cdot 2^{5-r} x^{5-r}$，令 $5-r=3$，得 $r=2$，可得 $x^3$ 项的系数为 $\mathrm{C}_5^2 \cdot (-1)^2 \cdot 2^{5-2} = 80$ .
\end{solution}
\end{question}

\begin{question}[points=4]
函数 $y = \cos x$ 在 $\left[-\frac{\pi}{2}, \frac{\pi}{4}\right]$ 上的值域为．
\fillin[{$[0,1]$}]
\begin{solution}
\textbf{答案：}$[0,1]$

\textbf{分析：}利用余弦函数的单调性可得.

\textbf{详解：}由函数 $y = \cos x$ 在 $\left[-\frac{\pi}{2}, 0\right]$ 上单调递增，在 $\left[0, \frac{\pi}{4}\right]$ 单调递减，且 $f(-\frac{\pi}{2}) = 0, f(0) = 1, f(\frac{\pi}{4}) = \frac{\sqrt{2}}{2}$，故函数 $y = \cos x$ 在 $\left[-\frac{\pi}{2}, \frac{\pi}{4}\right]$ 上的值域为 $[0,1]$ .
\end{solution}
\end{question}

\begin{question}[points=4]
已知随机变量 $X$ 的分布为 $\begin{pmatrix} 5 & 6 & 7 \\ 0.2 & 0.3 & 0.5 \end{pmatrix}$，则期望 $E[X] =$ ．
\fillin[{$6.3$}]
\begin{solution}
\textbf{答案：}$6.3$

\textbf{分析：}根据分布列结合期望公式可求期望.

\textbf{详解：}由题设有 $E[X] = 5 \times 0.2 + 6 \times 0.3 + 7 \times 0.5 = 1 + 1.8 + 3.5 = 6.3$ .
\end{solution}
\end{question}

\begin{question}[points=5]
如图，在正四棱柱 $ABCD-A_1B_1C_1D_1$ 中，$BD = 4\sqrt{2}, DB_1 = 9$，则该正四棱柱的体积为．

\begin{center}
  \includegraphics[width=.38\linewidth]{question-7-1.jpg}
\end{center}
\fillin[{$112$}]
\begin{solution}
\textbf{答案：}$112$

\textbf{分析：}求出侧棱长和底面边长后可求体积.

\textbf{详解：}因为 $BD = 4\sqrt{2}$ 且四边形 $ABCD$ 为正方形，故 $BA = 4$，而 $DB_1 = 9$，故 $BB_1^2 + BD^2 = 81$，故 $BB_1 = 7$，故所求体积为 $7 \times 16 = 112$ .
\end{solution}
\end{question}

\begin{question}[points=5]
设 $a, b > 0, a + \frac{1}{b} = 1$，则 $b + \frac{1}{a}$ 的最小值为．
\fillin[{$4$}]
\begin{solution}
\textbf{答案：}$4$

\textbf{分析：}灵活利用“1”将 $b + \frac{1}{a} = \left(b + \frac{1}{a}\right)\left(a + \frac{1}{b}\right)$ 展开利用基本不等式计算即可.

\textbf{详解：}易知 $b + \frac{1}{a} = \left(b + \frac{1}{a}\right)\left(a + \frac{1}{b}\right) = ab + \frac{1}{ab} + 2 \geq 2\sqrt{ab \cdot \frac{1}{ab}} + 2 = 4$，当且仅当 $ab = 1$，即 $a = \frac{1}{2}, b = 2$ 时取得最小值.
\end{solution}
\end{question}

\begin{question}[points=5]
4 个家长和 2 个儿童去爬山，6 个人需要排成一条队列，要求队列的头和尾均是家长，则不同的排列个数有种．
\fillin[{$288$}]
\begin{solution}
\textbf{答案：}$288$

\textbf{分析：}先选家长作队尾和队首，再排中间四人即可.

\textbf{详解：}先选两位家长排在首尾有 $\mathrm{P}_4^2 = 12$ 种排法；再排对中的四人有 $\mathrm{P}_4^4 = 24$ 种排法，故有 $12 \times 24 = 288$ 种排法.
\end{solution}
\end{question}

\begin{question}[points=5]
已知复数 $z$ 满足 $z^2 = (\overline{z})^2, |z| \leq 1$，则 $|z - 2 - 3i|$ 的最小值是．
\fillin[{$2\sqrt{2}$}]
\begin{solution}
\textbf{答案：}$2\sqrt{2}$

\textbf{分析：}先设 $z = a + bi$，利用复数的乘方运算及概念确定 $ab = 0$，再根据复数的几何意义数形结合计算即可.

\textbf{详解：}设 $z = a + bi\,(a,b \in \mathbf{R})$，则 $\overline{z} = a - bi$，由题意可知 $z^2 = a^2 + 2abi - b^2 = \overline{z}^2 = a^2 - 2abi - b^2$，则 $ab = 0$，又 $|z| = \sqrt{a^2 + b^2} \leq 1$，由复数的几何意义知 $z$ 在复平面内对应的点 $Z(a,b)$ 在单位圆内部（含边界）的坐标轴上运动，如图所示即线段 $AB, CD$ 上运动，设 $E(2,3)$，则 $|z-2-3i| = ZE$，由图象可知 $BE = \sqrt{10} > CE = 2\sqrt{2}$，所以 $ZE_{\min} = 2\sqrt{2}$ .
\end{solution}
\end{question}

\begin{question}[points=5]
小申同学观察发现，生活中有些时候影子可以完全投射在斜面上．某斜面上有两根长为 1 米的垂直于水平面放置的杆子，与斜面的接触点分别为 A、B，它们在阳光的照射下呈现出影子，阳光可视为平行光：其中一根杆子的影子在水平面上，长度为 0.4 米；另一根杆子的影子完全在斜面上，长度为 0.45 米．则斜面的底角 $\theta =$ ．（结果用角度制表示，精确到 $0.01^\circ$）
\fillin[{$12.58^\circ$}]
\begin{solution}
\textbf{答案：}$12.58^\circ$

\textbf{分析：}先根据在 A 处的旗杆算出阳光和水平面的夹角，然后结合 B 处的旗杆算出斜面角.

\textbf{详解：}如图，在 A 处，$\tan x = \frac{1}{0.4} = 2.5$，在 B 处满足 $\tan \angle CED = 2.5$，（其中 $ED //$ 水平面，$CE$ 是射过 B 处杆子最高点的光线，光线交斜面于 $E$），故设 $BD = y$，则 $ED = \frac{1+y}{2.5}$，由勾股定理，$y^2 + \left(\frac{1+y}{2.5}\right)^2 = 0.45^2$，解得 $y \approx 0.098$，于是 $\theta = \arcsin\frac{0.098}{0.45} \approx 12.58^\circ$ .
\end{solution}
\end{question}

\begin{question}[points=5]
已知 $f(x) = \begin{cases} 1, & x > 0 \\ 0, & x = 0 \\ -1, & x < 0 \end{cases}$，$\vec{a}, \vec{b}, \vec{c}$ 是平面内三个不同的单位向量．若 $f(\vec{a} \cdot \vec{b}) + f(\vec{b} \cdot \vec{c}) + f(\vec{c} \cdot \vec{a}) = 0$，则 $|\vec{a} + \vec{b} + \vec{c}|$ 的取值范围是．
\fillin[{$(1, \sqrt{5})$}]
\begin{solution}
\textbf{答案：}$(1, \sqrt{5})$

\textbf{分析：}利用分段函数值分类讨论，可得 $\{f(\vec{a} \cdot \vec{b}), f(\vec{b} \cdot \vec{c}), f(\vec{c} \cdot \vec{a})\} = \{-1, 0, 1\}$，再根据数量积关系设出 $\vec{a}, \vec{b}, \vec{c}$ 坐标，利用坐标运算，结合三角恒等变换求解模的范围可得.

\textbf{详解：}若 $f(\vec{a} \cdot \vec{b}) = f(\vec{b} \cdot \vec{c}) = f(\vec{c} \cdot \vec{a}) = 0$，则 $\vec{a} \cdot \vec{b} = \vec{b} \cdot \vec{c} = \vec{c} \cdot \vec{a} = 0$，又三个向量均为平面内的单位向量，故向量 $\vec{a}, \vec{b}, \vec{c}$ 两两垂直，显然不成立；故 $\{f(\vec{a} \cdot \vec{b}), f(\vec{b} \cdot \vec{c}), f(\vec{c} \cdot \vec{a})\} = \{-1, 0, 1\}$ . 不妨设 $\begin{cases} f(\vec{a} \cdot \vec{b}) = 1 \\ f(\vec{b} \cdot \vec{c}) = 0 \\ f(\vec{c} \cdot \vec{a}) = -1 \end{cases}$，则 $\vec{a} \cdot \vec{b} > 0, \vec{b} \cdot \vec{c} = 0, \vec{c} \cdot \vec{a} < 0$，不妨设 $\vec{b} = (1,0), \vec{c} = (0,1)$，$\vec{a} = (\cos \theta, \sin \theta), \theta \in [0, 2\pi)$，则 $\begin{cases} \vec{a} \cdot \vec{b} = \cos \theta > 0 \\ \vec{c} \cdot \vec{a} = \sin \theta < 0 \end{cases}$，则 $\theta \in \left(\frac{3\pi}{2}, 2\pi\right)$，则 $|\vec{a} + \vec{b} + \vec{c}| = \sqrt{(1+\cos \theta)^2 + (1+\sin \theta)^2} = \sqrt{3 + 2\cos \theta + 2\sin \theta} = \sqrt{3 + 2\sqrt{2}\sin\left(\theta + \frac{\pi}{4}\right)}$，由 $\theta \in \left(\frac{3\pi}{2}, 2\pi\right)$，$\theta + \frac{\pi}{4} \in \left(\frac{7\pi}{4}, \frac{9\pi}{4}\right)$，则 $\sin\left(\theta + \frac{\pi}{4}\right) \in \left(-\frac{\sqrt{2}}{2}, \frac{\sqrt{2}}{2}\right)$，$2\sqrt{2}\sin\left(\theta + \frac{\pi}{4}\right) \in (-2, 2)$，故 $|\vec{a} + \vec{b} + \vec{c}| \in (1, \sqrt{5})$ .
\end{solution}
\end{question}

\section*{二、选择题}
本大题共 4 题，第 13、14 题每题 4 分，第 15、16 题每题 5 分，共 18 分.每题有且仅有一个正确选项，考生应在答题纸的相应位置，将代表正确选项的小方格涂黑.

\begin{question}[points=4]
己知事件 $A$、$B$ 相互独立，事件 $A$ 发生的概率为 $P(A) = \frac{1}{2}$，事件 $B$ 发生的概率为 $P(B) = \frac{1}{2}$，则事件 $A \cap B$ 发生的概率 $P(A \cap B)$ 为（ ）
\begin{choices}
  \item $\frac{1}{8}$
  \item $\frac{1}{4}$
  \item $\frac{1}{2}$
  \item $0$
\end{choices}
\begin{solution}
\textbf{答案：}B

\textbf{分析：}根据独立事件的概率公式可求 $P(A \cap B)$ .

\textbf{详解：}因为 $A, B$ 相互独立，故 $P(A \cap B) = P(A)P(B) = \frac{1}{2} \times \frac{1}{2} = \frac{1}{4}$，故选：B.
\end{solution}
\end{question}

\begin{question}[points=4]
设 $a > 0, s \in \mathbf{R}$．下列各项中，能推出 $a^s > a$ 的一项是（ ）
\begin{choices}
  \item $a > 1$，且 $s > 0$
  \item $a > 1$，且 $s < 0$
  \item $0 < a < 1$，且 $s > 0$
  \item $0 < a < 1$，且 $s < 0$
\end{choices}
\begin{solution}
\textbf{答案：}D

\textbf{分析：}利用指数函数的性质分类讨论 $a$ 与 $1$ 的关系即可判定选项.

\textbf{详解：}$\because a > 0, a^s > a$，$\therefore a^{s-1} > 1 = a^0$，当 $a \in (0,1)$ 时，$y = a^x$ 定义域上严格单调递减，此时若 $s-1 < 0$，则一定有 $a^{s-1} > 1 = a^0$ 成立，故 D 正确，C 错误；当 $a \in (1, +\infty)$ 时，$y = a^x$ 定义域上严格单调递增，要满足 $a^{s-1} > 1 = a^0$，需 $s > 1$，即 A、B 错误. 故选：D.
\end{solution}
\end{question}

\begin{question}[points=5]
已知 $A(0,1), B(1,2)$，$C$ 在 $\Gamma: x^2 - y^2 = 1 (x \geq 1, y \geq 0)$ 上，则 $\triangle ABC$ 的面积（ ）
\begin{choices}
  \item 有最大值，但没有最小值
  \item 没有最大值，但有最小值
  \item 既有最大值，也有最小值
  \item 既没有最大值，也没有最小值
\end{choices}
\begin{solution}
\textbf{答案：}A

\textbf{分析：}设出曲线上一点为 $(a, b)$，得出 $a = \sqrt{b^2 + 1}$，将三角形的高转化成关于 $b$ 的函数，分析其单调性，从而求解.

\textbf{详解：}设曲线上一点为 $(a, b)$，则 $a^2 - b^2 = 1$，则 $a = \sqrt{b^2 + 1}$，$k_{AB} = \frac{2-1}{1-0} = 1$，$AB$ 方程为：$y-1 = x$，即 $x - y + 1 = 0$，根据点到直线的距离公式，$(a, b)$ 到 $AB$ 的距离为：$\frac{|a - b + 1|}{\sqrt{2}} = \frac{\sqrt{b^2 + 1} - b + 1}{\sqrt{2}}$，设 $f(b) = \sqrt{b^2 + 1} - b = \frac{1}{\sqrt{b^2 + 1} + b}$，由于 $b \geq 0$，显然 $f(b)$ 关于 $b$ 单调递减，$f(b)_{\max} = f(0)$，无最小值，即 $\triangle ABC$ 中，$AB$ 边上的高有最大值，无最小值，又 $AB$ 一定，故面积有最大值，无最小值. 故选：A.
\end{solution}
\end{question}

\begin{question}[points=5]
已知数列 $\{a_n\}$、$\{b_n\}$、$\{c_n\}$ 的通项公式分别为 $a_n = 10n - 9$，$b_n = 2^n$，$c_n = \lambda a_n + (1-\lambda)b_n$．若对任意的 $\lambda \in [0,1]$，$a_n$、$b_n$、$c_n$ 的值均能构成三角形，则满足条件的正整数 $n$ 有（ ）
\begin{choices}
  \item $4$ 个
  \item $3$ 个
  \item $1$ 个
  \item 无数个
\end{choices}
\begin{solution}
\textbf{答案：}B

\textbf{分析：}由 $c_n = \lambda a_n + (1-\lambda)b_n$ 可知 $c_n$ 范围，再由三角形三边关系可得 $a_n, b_n, c_n$ 的不等关系，结合函数零点解不等式可得.

\textbf{详解：}由题意 $a_n, b_n, c_n > 0$，不妨设 $A(n, a_n), B(n, b_n), C(n, c_n)$，三点均在第一象限内，由 $c_n = \lambda a_n + (1-\lambda)b_n$ 可知，$\overrightarrow{BC} = \lambda \overrightarrow{BA}, \lambda \in [0,1]$，故点 $C$ 恒在线段 $AB$ 上，则有 $\min\{a_n, b_n\} \leq c_n \leq \max\{a_n, b_n\} < a_n + b_n$．即对任意的 $\lambda \in [0,1]$，$c_n < a_n + b_n$ 恒成立，令 $10x - 9 = 2^x$，构造函数 $f(x) = 2^x - 10x + 9, x > 0$，则 $f'(x) = 2^x \ln 2 - 10$，由 $f'(x)$ 单调递增，又 $f'(3) < 0, f'(4) > 0$，存在 $x_0 \in (3,4)$，使 $f'(x_0) = 0$，即当 $0 < x < x_0$ 时，$f'(x) < 0$，$f(x)$ 单调递减；当 $x > x_0$ 时，$f'(x) > 0$，$f(x)$ 单调递增；故 $f(x)$ 至多 $2$ 个零点，又由 $f(1) > 0, f(2) < 0, f(5) < 0, f(6) > 0$，可知 $f(x)$ 存在 $2$ 个零点，不妨设 $x_1, x_2 (x_1 < x_2)$，且 $x_1 \in (1,2), x_2 \in (5,6)$．①若 $a_n \leq b_n$，即 $10n - 9 \leq 2^n$ 时，此时 $n = 1$ 或 $n \geq 6$．则 $a_n \leq c_n \leq b_n$，可知 $b_n + c_n > a_n$ 成立，要使 $a_n$、$b_n$、$c_n$ 的值均能构成三角形，所以 $a_n + c_n > b_n$ 恒成立，故 $b_n < 2a_n$，所以有 $\begin{cases} 10n - 9 \leq 2^n \\ 2^n < 2(10n - 9) \end{cases}$，解得 $n = 6$；②若 $a_n \geq b_n$，即 $10n - 9 \geq 2^n$ 时，此时 $n = 2,3,4,5$．则 $a_n \geq c_n \geq b_n$，可知 $a_n + c_n > b_n$ 成立，要使 $a_n$、$b_n$、$c_n$ 的值均能构成三角形，所以 $b_n + c_n > a_n$ 恒成立，故 $a_n < 2b_n$，所以有 $\begin{cases} 10n - 9 \geq 2^n \\ 10n - 9 < 2^{n+1} \end{cases}$，解得 $n = 4$ 或 $5$；综上可知，正整数 $n$ 的个数有 $3$ 个. 故选：B.
\end{solution}
\end{question}

\section*{三、解答题}
本大题共 5 题，第 17-19 题每题 14 分，第 20-21 题每题 18 分，共 78 分.解答下列各题必须在答题纸的相应位置写出必要的步骤.

\begin{question}[points=14]
2024 年东京奥运会，中国获得了男子 $4 \times 100$ 米混合泳接力金牌．以下是历届奥运会男子 $4 \times 100$ 米混合泳接力项目冠军成绩记录（单位：秒），数据按照升序排列．



206.78  207.46  207.95  209.34  209.35

210.68  213.73  214.84  216.93  216.93



(1) 求这组数据的极差与中位数；

(2) 从这 10 个数据中任选 3 个，求恰有 2 个数据在 211 以上的概率；

(3) 若比赛成绩 $y$ 关于年份 $x$ 的回归方程为 $\hat{y} = -0.311x + \hat{b}$，年份 $x$ 的平均数为 $2006$，预测 $2028$ 年冠军队的成绩（精确到 $0.01$ 秒）．
\begin{solution}
\textbf{答案：}(1) 极差 $10.15$，中位数 $210.015$； (2) $\frac{3}{10}$； (3) $204.56$ 秒

\textbf{分析：}(1) 由最长与最短用时可得极差，由中间两数平均数可得中位数； (2) 由古典概型概率公式可得； (3) 先求成绩平均数 $y$，再由 $(\bar{x}, \bar{y})$ 在回归直线上，代入方程可得 $\hat{b}$，再代入年份预测可得.

\textbf{详解：}(1) 由题意，数据的最大值为 $216.93$，最小值为 $206.78$，则极差为 $216.93 - 206.78 = 10.15$；数据中间两数为 $209.35$ 与 $210.68$，则中位数为 $\frac{209.35 + 210.68}{2} = 210.015$．故极差为 $10.15$，中位数为 $210.015$．
(2) 由题意，数据共 $10$ 个，$211$ 以上数据共有 $4$ 个，故设事件 $A = $“恰有 $2$ 个数据在 $211$ 以上”，则 $P(A) = \frac{\mathrm{C}_4^2 \cdot \mathrm{C}_6^1}{\mathrm{C}_{10}^3} = \frac{3}{10}$．故恰有 $2$ 个数据在 $211$ 以上的概率为 $\frac{3}{10}$．
(3) 由题意，成绩的平均数 $\bar{y} = \frac{206.78 + 207.46 + 207.95 + 209.34 + 209.35 + 210.68 + 213.73 + 214.84 + 216.93 + 216.93}{10} = 211.399$，由直线 $\hat{y} = -0.311x + \hat{b}$ 过 $(2006, 211.399)$，则 $\hat{b} = 211.399 + 0.311 \times 2006 = 835.265$，故回归直线方程为 $\hat{y} = -0.311x + 835.265$．当 $x = 2028$ 时，$\hat{y} = -0.311 \times 2028 + 835.265 = 204.557 \approx 204.56$．故预测 $2028$ 年冠军队的成绩为 $204.56$ 秒.
\end{solution}
\end{question}

\begin{question}[points=14]
如图，$P$ 是圆锥的顶点，$O$ 是底面圆心，$AB$ 是底面直径，且 $AB = 2$．



(1) 若直线 $PA$ 与圆锥底面的所成角为 $\frac{\pi}{3}$，求圆锥的侧面积；

(2) 已知 $Q$ 是母线 $PA$ 的中点，点 $C$、$D$ 在底面圆周上，且弧 $AC$ 的长为 $\frac{\pi}{3}$，$CD \parallel AB$．设点 $M$ 在线段 $OC$ 上，证明：直线 $QM \parallel$ 平面 $PBD$．

\begin{center}
  \includegraphics[width=.38\linewidth]{question-18-1.jpg}
\end{center}
\begin{solution}
\textbf{答案：}(1) $2\pi$； (2) 证明见解析

\textbf{分析：}(1) 由线面角先算出母线长，然后根据侧面积公式求解； (2) 证明平面 $QOC \parallel$ 平面 $PBD$，然后根据面面平行的性质可得.

\textbf{详解：}(1) 由题知，$\angle PAB = \frac{\pi}{3}$，即轴截面 $\triangle ABP$ 是等边三角形，故 $PA = AB = 2$，底面周长为 $2\pi \times 1 = 2\pi$，则侧面积为：$\frac{1}{2} \times 2 \times 2\pi = 2\pi$．
(2) 由题知 $AQ = QP, AO = OB$，则根据中位线性质，$QO \parallel PB$，又 $QO \not\subset$ 平面 $PBD$，$PB \subset$ 平面 $PBD$，则 $QO \parallel$ 平面 $PBD$．由于 $\widehat{AC} = \frac{\pi}{3}$，底面圆半径是 $1$，则 $\angle AOC = \frac{\pi}{3}$，又 $CD \parallel AB$，则 $\angle OCD = \frac{\pi}{3}$，又 $OC = OD$，则 $\triangle OCD$ 为等边三角形，则 $CD = 1$，于是 $CD \parallel BO$ 且 $CD = OB$，则四边形 $OBDC$ 是平行四边形，故 $OC \parallel BD$，又 $OC \not\subset$ 平面 $PBD$，$BD \subset$ 平面 $PBD$，故 $OC \parallel$ 平面 $PBD$．又 $OC \cap OQ = O$，$OC, OQ \subset$ 平面 $QOC$，根据面面平行的判定，于是平面 $QOC \parallel$ 平面 $PBD$，又 $M \in OC$，则 $QM \subset$ 平面 $QOC$，则 $QM \parallel$ 平面 $PBD$．
\end{solution}
\end{question}

\begin{question}[points=14]
已知 $f(x) = x^2 - (m+2)x + m\ln x, m \in \mathbf{R}$．

(1) 若 $f(1) = 0$，求不等式 $f(x) \leq x^2 - 1$ 的解集；

(2) 若函数 $y = f(x)$ 满足在 $(0, +\infty)$ 上存在极大值，求 $m$ 的取值范围；
\begin{solution}
\textbf{答案：}(1) $[1, +\infty)$； (2) $m > 0$ 且 $m \neq 2$

\textbf{分析：}(1) 先求出 $m$，从而原不等式即为 $x + \ln x \geq 1$，构建新函数 $s(x) = x + \ln x, x > 0$，由该函数为增函数可求不等式的解； (2) 求出函数的导数，就 $m \leq 0, 0 < m < 2, m = 2, m > 2$ 分类讨论后可得参数的取值范围.

\textbf{详解：}(1) 因为 $f(1) = 0$，故 $1 - m - 2 + 0 = 0$，故 $m = -1$，故 $f(x) = x^2 - x - \ln x$，故 $f(x) \leq x^2 - 1$ 即为 $x + \ln x \geq 1$，设 $s(x) = x + \ln x, x > 0$，则 $s'(x) = 1 + \frac{1}{x} > 0$，故 $s(x)$ 在 $(0, +\infty)$ 上为增函数，而 $x + \ln x \geq 1$ 即为 $s(x) \geq s(1)$，故 $x \geq 1$，故原不等式的解为 $[1, +\infty)$．
(2) $f(x)$ 在 $(0, +\infty)$ 有极大值即为有极大值点．$f'(x) = 2x - (m+2) + \frac{m}{x} = \frac{2x^2 - (m+2)x + m}{x} = \frac{(2x - m)(x - 1)}{x}$．若 $m \leq 0$，则 $x \in (0,1)$ 时，$f'(x) < 0$，$x \in (1, +\infty)$ 时，$f'(x) > 0$，故 $x = 1$ 为 $f(x)$ 的极小值点，无极大值点，故舍；若 $0 < \frac{m}{2} < 1$ 即 $0 < m < 2$，则 $x \in \left(\frac{m}{2}, 1\right)$ 时，$f'(x) < 0$，$x \in \left(0, \frac{m}{2}\right) \cup (1, +\infty)$ 时，$f'(x) > 0$，故 $x = \frac{m}{2}$ 为 $f(x)$ 的极大值点，符合题设要求；若 $m = 2$，则 $x \in (0, +\infty)$ 时，$f'(x) \geq 0$，$f(x)$ 无极值点，舍；若 $\frac{m}{2} > 1$ 即 $m > 2$，则 $x \in \left(1, \frac{m}{2}\right)$ 时，$f'(x) < 0$，$x \in (0,1) \cup \left(\frac{m}{2}, +\infty\right)$ 时，$f'(x) > 0$，故 $x = 1$ 为 $f(x)$ 的极大值点，符合题设要求；综上，$m > 0$ 且 $m \neq 2$.
\end{solution}
\end{question}

\begin{question}[points=18]
已知椭圆 $\Gamma: \frac{x^2}{a^2} + \frac{y^2}{5} = 1 (a > \sqrt{5})$，$M(0, m) (m > 0)$，$A$ 是 $\Gamma$ 的右顶点．

(1) 若 $\Gamma$ 的焦点 $(2,0)$，求离心率 $e$；

(2) 若 $a = 4$，且 $\Gamma$ 上存在一点 $P$，满足 $\overrightarrow{PA} = 2\overrightarrow{MP}$，求 $m$；

(3) 已知 $AM$ 的中垂线 $l$ 的斜率为 $2$，$l$ 与 $\Gamma$ 交于 $C$、$D$ 两点，$\angle CMD$ 为钝角，求 $a$ 的取值范围．
\begin{solution}
\textbf{答案：}(1) $\frac{2}{3}$； (2) $\sqrt{10}$； (3) $(\sqrt{5}, \sqrt{11})$

\textbf{分析：}(1) 由方程可得 $b^2 = 5$，再由焦点坐标得 $c$，从而求出 $a$ 得离心率； (2) 设点 $P$ 坐标，由向量关系 $\overrightarrow{PA} = 2\overrightarrow{MP}$ 坐标化可解得 $P$ 坐标，代入椭圆方程可得 $m$； (3) 根据中垂线性质，由斜率与中点坐标得直线 $l$ 方程，联立直线与椭圆方程，将钝角条件转化为向量不等式 $\overrightarrow{MC} \cdot \overrightarrow{MD} < 0$，再坐标化利用韦达定理代入化简不等式求解可得 $a$ 范围.

\textbf{详解：}(1) 由题意知，$\Gamma: \frac{x^2}{a^2} + \frac{y^2}{5} = 1 (a > \sqrt{5})$，则 $b^2 = 5$，由右焦点 $(2,0)$，可知 $c = 2$，则 $a = \sqrt{5 + c^2} = 3$，故离心率 $e = \frac{c}{a} = \frac{2}{3}$．
(2) 由题意 $A(4,0), M(0, m) (m > 0)$，$P(x_P, y_P)$，由 $\overrightarrow{PA} = 2\overrightarrow{MP}$ 得，$\begin{cases} 4 - x_P = 2x_P \\ -y_P = 2y_P - 2m \end{cases}$，解得 $P\left(\frac{4}{3}, \frac{2m}{3}\right)$，代入 $\frac{x^2}{16} + \frac{y^2}{5} = 1$，得 $\frac{1}{9} + \frac{4m^2}{45} = 1$，又 $m > 0$，解得 $m = \sqrt{10}$．
(3) 由线段 $AM$ 的中垂线 $l$ 的斜率为 $2$，所以直线 $AM$ 的斜率为 $-\frac{1}{2}$，则 $\frac{m-0}{0-a} = -\frac{1}{2}$，解得 $m = \frac{a}{2}$，由 $A(a,0), M(0, \frac{a}{2})$ 得 $AM$ 中点坐标为 $(\frac{a}{2}, \frac{a}{4})$，故直线 $l: y = 2x - \frac{3}{4}a$，显然直线 $l$ 过椭圆内点 $(\frac{3}{8}a, 0)$，故直线与椭圆恒有两不同交点，设 $C(x_1, y_1), D(x_2, y_2)$，由 $\begin{cases} y = 2x - \frac{3}{4}a \\ 5x^2 + a^2 y^2 = 5a^2 \end{cases}$ 消 $y$ 得 $(4a^2 + 5)x^2 - 3a^3 x + \frac{9}{16}a^4 - 5a^2 = 0$，由韦达定理得 $x_1 + x_2 = \frac{3a^3}{4a^2 + 5}$，$x_1 x_2 = \frac{\frac{9}{16}a^4 - 5a^2}{4a^2 + 5}$．因为 $\angle CMD$ 为钝角，则 $\overrightarrow{MC} \cdot \overrightarrow{MD} < 0$，且 $M(0, \frac{a}{2})$，则有 $x_1 x_2 + (y_1 - \frac{a}{2})(y_2 - \frac{a}{2}) < 0$，所以 $x_1 x_2 + \left(2x_1 - \frac{5a}{4}\right)\left(2x_2 - \frac{5a}{4}\right) = 5x_1 x_2 - \frac{5a}{2}(x_1 + x_2) + \frac{25}{16}a^2 < 0$，即 $5\left(\frac{\frac{9}{16}a^4 - 5a^2}{4a^2 + 5}\right) - \frac{5a}{2} \cdot \frac{3a^3}{4a^2 + 5} + \frac{25}{16}a^2 < 0$，解得 $a^2 < 11$，又 $a > \sqrt{5}$，故 $\sqrt{5} < a < \sqrt{11}$，即 $a$ 的取值范围是 $(\sqrt{5}, \sqrt{11})$．
\end{solution}
\end{question}

\begin{question}[points=18]
已知函数 $y = f(x)$ 的定义域为 $\mathbf{R}$．对于正实数 $a$，定义集合 $M_a = \{x \mid f(x+a) = f(x)\}$．

(1) 若 $f(x) = \sin x$，判断 $\frac{\pi}{3}$ 是否是 $M_\pi$ 中的元素，请说明理由；

(2) 若 $f(x) = \begin{cases} x+2, & x < 0 \\ x, & x \geq 0 \end{cases}$，$M_a \neq \varnothing$，求 $a$ 的取值范围；

(3) 若 $y = f(x)$ 是偶函数，当 $x \in (0,1]$ 时，$f(x) = 1 - x$，且对任意 $a \in (0,2)$，均有 $M_a \subseteq M_2$．写出 $y = f(x)$，$x \in (1,2)$ 解析式，并证明：对任意实数 $c$，函数 $y = f(x) - c$ 在 $[-3,3]$ 上至多有 $9$ 个零点．
\begin{solution}
\textbf{答案：}(1) 不是； (2) $\left[\frac{7}{4}, 4\right)$； (3) 解析式 $f(x) = x - 1, x \in (1,2)$；证明见解析

\textbf{分析：}(1) 直接代入计算 $f\left(\frac{\pi}{3}\right)$ 和 $f\left(\frac{\pi}{3} + \pi\right)$ 即可； (2) 法一：转化为存在实数 $x_0$ 使得 $f(x_0 + a) = f(x_0)$，分析得 $x_0 + 2 = \sqrt{x_0 + a}$，再计算得 $a = (x_0 + 2)^2 - x_0 = \left(x_0 + \frac{3}{2}\right)^2 + \frac{7}{4}$，最后根据 $x_0$ 的范围即可得到答案；法二：画出函数图象，转化为直线 $y = t$ 与该函数有两个交点，将 $a$ 用 $t$ 表示，最后利用二次函数性质即可得到答案； (3) 利用函数奇偶性和集合新定义即可求出 $x \in (1,2)$ 时解析式，再分析出 $f(-3) \notin (0,1)$，最后对 $c$ 的范围进行分类讨论即可.

\textbf{详解：}(1) $f\left(\frac{\pi}{3}\right) = \sin\frac{\pi}{3} = \frac{\sqrt{3}}{2}$，$f\left(\frac{\pi}{3} + \pi\right) = -\sin\frac{\pi}{3} = -\frac{\sqrt{3}}{2}$，两者不相等，故 $\frac{\pi}{3}$ 不是 $M_\pi$ 中的元素．
(2) 法一：因为 $M_a \neq \varnothing$，则存在实数 $x_0$ 使得 $f(x_0 + a) = f(x_0)$，且 $a > 0$，当 $x < 0$ 时，$f(x) = x+2$，其在 $(-\infty, 0)$ 上严格单调递增，当 $x \geq 0$ 时，$f(x) = \sqrt{x}$，其在 $[0, +\infty)$ 上也严格单调递增，则 $x_0 < 0 \leq x_0 + a$，则 $x_0 + 2 = \sqrt{x_0 + a}$，令 $x+2 = 0$，解得 $x = -2$，则 $-2 \leq x_0 < 0$，则 $a = (x_0 + 2)^2 - x_0 = \left(x_0 + \frac{3}{2}\right)^2 + \frac{7}{4} \in \left[\frac{7}{4}, 4\right)$．法二：作出该函数图象，则由题意知直线 $y = t$ 与该函数有两个交点，由图知 $0 \leq t < 2$，假设交点分别为 $A(m, t)$，$B(n, t)$，联立方程组 $\begin{cases} m = t^2 \\ n+2 = t \end{cases}$ 得 $a = |AB| = m - n = t^2 - (t-2) = \left(t - \frac{1}{2}\right)^2 + \frac{7}{4} \in \left[\frac{7}{4}, 4\right)$．
(3) 对任意 $x_0 \in (1,2)$，$x_0 - 2 \in (-1,0)$，因为其是偶函数，则 $f(x_0 - 2) = f(2 - x_0)$，而 $4 - 2x_0 \in (0,2)$，所以 $x_0 - 2 \in M_{4-2x_0} \subseteq M_2$，所以 $f(x_0) = f(x_0 - 2) = f(2 - x_0)$，因为 $x_0 \in (1,2)$，则 $2 - x_0 \in (0,1)$，所以 $f(x_0) = f(2 - x_0) = 1 - (2 - x_0) = x_0 - 1$，所以 $f(x) = x - 1, x \in (1,2)$．所以当 $s \in (0,1)$ 时，$1-s \in (0,1)$，$1+s \in (1,2)$，则 $f(1-s) = 1 - (1-s) = s$，$f(1+s) = (1+s) - 1 = s$，则 $f(1-s) = f(1+s)$，而 $1+s - (1-s) = 2s$，$(3-s) - (1-s) = 2$，则 $1-s \in M_{2s} \subseteq M_2$，则 $f(1-s) = f(3-s)$，所以当 $x \in (2,3)$ 时，$f(x) = f(x-2) = 1 - (x-2) = 3 - x$，而 $f(x)$ 为偶函数，画出函数图象如下：其中 $f(-3) = f(3)$，$f(-2) = f(2)$，$f(0)$，但其对应的 $y$ 值均未知．首先说明 $f(-3) = n \notin (0,1)$，若 $f(-3) = n \in (0,1)$，则 $-3+n \in (-3,-2)$，易知此时 $f(x) = x+3, x \in (-3,-2)$，则 $f(-3+n) = n$，所以 $f(-3) \in M_n \subseteq M_2$，而 $x \in [-1,0)$ 时，$f(x) = x+1$，所以 $f(-3) = f(-1) = 0$，与 $f(-3) = n$ 矛盾，所以 $f(-3) \notin (0,1)$，即 $f(-3) = f(3) \notin (0,1)$．令 $y = f(x) - c = 0$，则 $y = f(x) = c$，当 $c = 0$ 时，即使让 $f(-3) = f(3) = f(-2) = f(2) = f(0) = 0$，此时最多 $7$ 个零点；当 $c \geq 1$ 时，若 $f(-2) = f(2) = f(0) = f(-3) = f(3) = c$，此时有 $5$ 个零点，故此时最多 $5$ 个零点；当 $c < 0$ 时，若 $f(-2) = f(2) = f(0) = f(-3) = f(3) = c$，此时有 $5$ 个零点，故此时最多 $5$ 个零点；当 $0 < c < 1$ 时，若 $f(-2) = f(2) = f(0) = c$，此时有 $3$ 个零点，若 $f(-3) = c \in (0,1)$，则 $-3+c \in (-3,-2)$，易知此时 $f(x) = x+3$，则 $f(-3+c) = c$，所以 $f(-3) \in M_c \subseteq M_2$，而 $x \in [-1,0)$ 时，$f(x) = x+1$，所以 $f(-3) = f(-1) = 0$，与 $f(-3) = c$ 矛盾，所以 $f(-3) \notin (0,1)$，则最多在 $(-3,-2), (-2,-1), (-1,0), (0,1), (1,2), (2,3)$ 之间取得 $6$ 个零点，以及在 $x = -2, 0, 2$ 处成为零点，故不超过 $9$ 个零点．综上，零点不超过 $9$ 个.
\end{solution}
\end{question}

